\documentclass[crop=false]{standalone}
\usepackage{15150toolbox}
\usepackage{needspace}

% Toggle: set \soltrue to render solutions inline (blue boxes).
\newif\ifsol
\soltrue   % comment this line for the student version

% Solution environment. When \ifsol is true, render as a blue framed box;
% when false, use the comment package to swallow the body (handles verbatim
% content correctly, unlike environ's \BODY).
\usepackage{comment}
\ifsol
  \newenvironment{solution}
    {\par\color{solution_color}\begin{framed}\textbf{Solution.}\par}
    {\end{framed}}
\else
  \excludecomment{solution}
\fi

% Lecture-slide macros that aren't in 15150toolbox; local re-definitions
% so the quiz is self-contained.
\newcommand{\term}[1]{\textbf{\textit{#1}}}
\newcommand{\thmBox}[2]{\par\noindent\begin{mdframed}[backgroundcolor=blue!8,hidealllines=true]\textbf{Theorem#1.} #2\end{mdframed}}
\newcommand{\bcBox}[1]{\par\noindent\textbf{Base case.} #1\par}
\newcommand{\ihBox}[2]{\par\noindent\textbf{Inductive hypothesis#1.} #2\par}
\newcommand{\isBox}[2]{\par\noindent\textbf{Inductive step#1.} #2\par}

\newcommand{\quizproblem}[2]{%
  \needspace{5\baselineskip}%
  \stepcounter{problem}%
  \ifsol
    \subsection*{Problem \arabic{problem}. #1 \hfill \timing{#2}}%
  \else
    \subsection*{Problem \arabic{problem}. #1}%
  \fi
}

\newcommand{\answerspace}[1]{\ifsol\else\vspace{#1}\par\fi}

\printout{Quiz 1 --- Lectures 1--4}

\begin{document}

\begin{center}
  {\Large \textbf{15-150 --- Quiz 1}} \\
  \vspace{4pt}
  Topics: \textit{Prologue}; \textit{Equivalence, Binding, and Scope}; \\
  \textit{Induction and Recursion}; \textit{Structural Induction and Tail Recursion} \\
  \vspace{4pt}
  \textit{60 minutes. You should comfortably finish in 30--45 minutes.}
\end{center}

\vspace{1em}

\noindent\textbf{Name:} \hrulefill \hspace{2em} \textbf{Andrew ID:} \hrulefill

\vspace{1em}

% =============================================================================
% PROBLEM 1 --- Equivalence and Specifications
% =============================================================================

\quizproblem{Equivalence and Specifications}{10}

\textbf{(a)} A student writes the following definition of extensional
equivalence:

\begin{quote}
``Two expressions \code{e1} and \code{e2} of the same type are
\term{extensionally equivalent} if they evaluate to the same value.''
\end{quote}

This definition is \textbf{incomplete}. State what is missing, and give an
example of two expressions of the same type that the definition above
\emph{incorrectly} classifies as not equivalent.

\answerspace{4cm}

\begin{solution}
The definition is missing the cases where \code{e1} and \code{e2}
\textbf{do not evaluate to a value}. The full definition requires three
cases: they reduce to the same value, \emph{or} they both loop forever,
\emph{or} they both raise the same exception.

Acceptable examples: any two nonterminating expressions of the same
type, or any two expressions that raise the same exception.
\end{solution}

\vspace{1em}

\textbf{(b)} Consider the \code{pow} function from lecture, with a
partial specification:

\begin{codeblock}
  (* pow : int * int -> int *)
  (* REQUIRES: <fill this in> *)
  (* ENSURES: pow (n, k) evaluates to n^k. *)
  fun pow (n : int, 0 : int) : int = 1
    | pow (n, k) = n * pow (n, k - 1)
\end{codeblock}

\textbf{Write a correct REQUIRES clause for \code{pow}.}

\vspace{0.5em}

\hspace{2em} REQUIRES: \underline{\hspace{8cm}}

\answerspace{1.5cm}

\begin{solution}
\textbf{REQUIRES: \code{k >= 0}.}

Acceptable answers: \code{k >= 0}, \code{0 <= k}, ``\code{k} is
non-negative.'' A student who writes \code{n >= 0} or \code{n >= 0 and
k >= 0} has not analyzed which variable controls termination.
\end{solution}

% =============================================================================
% PROBLEM 2 --- Well-Typed vs. Valuable
% =============================================================================

\clearpage
\quizproblem{Well-Typed vs.\ Valuable}{5}

For each expression below, fill in both columns. For \textbf{Well-typed?}
write yes if the expression is well-typed (and give its type), or write no. For \textbf{Valuable?}
write yes if the expression reduces to a value (and give the value), or write no.

\vspace{0.5em}

\begin{center}
\begin{tabular}{|l|p{4cm}|p{4cm}|}
  \hline
  \textbf{Expression} & \textbf{Well-typed?} & \textbf{Valuable?} \\
  \hline
  \code{2 + 2} & & \\[2.2em] \hline
  \code{1 div 0} & & \\[2.2em] \hline
  \code{"hi" + 1} & & \\[2.2em] \hline
  \code{(fn x => x) 5} & & \\[2.2em] \hline
\end{tabular}
\end{center}

\begin{solution}
\begin{center}
\begin{tabular}{|l|l|l|}
  \hline
  \textbf{Expression} & \textbf{Well-typed?} & \textbf{Valuable?} \\
  \hline
  \code{2 + 2} & Yes, \code{int} & Yes, \code{4} \\
  \hline
  \code{1 div 0} & Yes, \code{int} & No, raises \code{Div} \\
  \hline
  \code{"hi" + 1} & No (\code{+} needs \code{int}s) & No \\
  \hline
  \code{(fn x => x) 5} & Yes, \code{int} & Yes, \code{5} \\
  \hline
\end{tabular}
\end{center}
\end{solution}

% =============================================================================
% PROBLEM 3 --- Binding and Scope
% =============================================================================

\quizproblem{Binding and Scope}{5}

Consider the following SML code:

\begin{codeblock}
  val x : int = 2
  fun foo (y : int) : int = x + y
  val x : int = 4
\end{codeblock}

After these three declarations have executed, the function \code{foo} is
\emph{later} called as \code{foo 1}.

\vspace{0.5em}

\textbf{(a)} Which environment does the body of \code{foo} use to look
up \code{x} during this call? Circle one.

\vspace{0.5em}

\hspace{2em}
(A) \quad $[\code{2}/\code{x}]$ \hfill
(B) \quad $[\code{4}/\code{x}]$ \hfill
(C) \quad $[\code{2}/\code{x},\ \code{4}/\code{x}]$ \hfill
(D) \quad whichever was bound most recently when \code{foo} was called

\vspace{1em}

\textbf{(b)} In one sentence, explain why.

\answerspace{4cm}

\begin{solution}
\textbf{(a)} \textbf{A.} The body of \code{foo} looks up \code{x} in
$[\code{2}/\code{x}]$.

\textbf{(b)} When \code{foo} was \emph{declared}, the binding
of \code{x} to \code{2} was active, so \code{foo} ``captures'' that binding and
carries it forward. Later bindings of \code{x} (the \code{val x = 4})
merely shadow the old one.

Would also accept any answer involving the "closure" of \code{foo}, or that
\code{foo} "saves" or "remembers" the environment at the time of declaration.
\end{solution}

% =============================================================================
% PROBLEM 4 --- Induction Proof
% =============================================================================

\quizproblem{An Induction Proof}{10}

Recall the \code{@} function from lecture:

\begin{codeblock}
  infix @
  fun ([] : int list) @ (R : int list) : int list = R
    | (x::xs) @ R = x :: (xs @ R)
\end{codeblock}

\thmBox{}{\, For all \code{L : int list}, \code{L @ []} $\eeq$ \code{L}.}

\textbf{Prove this theorem.} Be sure to state your inductive hypothesis
explicitly and justify each step of your calculation.

\answerspace{11cm}

\begin{solution}
We proceed by \textbf{structural induction} on \code{L : int list}.

\vspace{5pt}

\bcBox{\, Case: \code{L = []}}
\begin{align*}
  \code{[] @ []}  &\stepsTo \code{[]}    & \text{(clause 1 of \code{@})}
\end{align*}

\vspace{5pt}

\ihBox{}{\, Case: \code{L = xs}, for some \code{xs : int list}. Assume
\code{xs @ []} $\eeq$ \code{xs}.}

\vspace{5pt}

\isBox{}{\, Case: \code{L = x::xs}, for some \code{x : int}. We want to
show \code{(x::xs) @ []} $\eeq$ \code{x::xs}.}

\begin{align*}
  \code{(x::xs) @ []}
    &\stepsTo \code{x :: (xs @ [])}
        & \text{(clause 2 of \code{@})} \\
    &\eeq    \code{x :: xs}
        & \text{(inductive hypothesis)}
\end{align*}

Thus, by the principle of structural induction, the theorem holds for
all \code{L : int list}.

\vspace{5pt}

\textbf{Common mistakes to watch for:}
\begin{itemize}
  \item Using simple/natural-number induction instead of structural ---
    \code{L} is a list, not a number.
  \item Quantifying the IH wrong (``Assume for \emph{all} \code{xs},
    \code{xs @ []} $\eeq$ \code{xs}'') --- that assumes the theorem.
    The IH is for \emph{some} \code{xs}.
\end{itemize}
\end{solution}

% =============================================================================
% PROBLEM 5 --- Writing a Tail-Recursive Function
% =============================================================================

\clearpage
\quizproblem{A Tail-Recursive \code{sum}}{10}

In lecture, we wrote a non-tail-recursive \code{length} function, and then
rewrote it tail-recursively using an \term{accumulator}:

\begin{codeblock}
  fun tlength ([] : int list, acc : int) : int = acc
    | tlength (x::xs, acc) = tlength (xs, 1 + acc)

  fun length (L : int list) : int = tlength (L, 0)
\end{codeblock}

Following the same pattern, write a tail-recursive function \code{tsum}
and a wrapper \code{sum} satisfying the spec:

\spec
  {sum}
  {int list -> int}
  {true}
  {\code{sum L} evaluates to the sum of the elements of \code{L}.}

\answerspace{4cm}

\begin{solution}
\begin{codeblock}
  fun tsum ([] : int list, acc : int) : int = acc
    | tsum (x::xs, acc) = tsum (xs, x + acc)

  fun sum (L : int list) : int = tsum (L, 0)
\end{codeblock}
\end{solution}

% =============================================================================
% PROBLEM 6 --- Reduction vs. Equivalence
% =============================================================================

\quizproblem{Reduction vs. Equivalence}{5}

Recall the \code{fast\_pow} and \code{pow} functions from lecture. The
following is an attempted proof of the base case of the theorem
\code{fast\_pow} $\eeq$ \code{pow}:

\begin{align*}
  \code{fast\_pow (n, 0)}
    &\stepsTo \code{1}              & \text{(clause 1 of \code{fast\_pow})} \\
    &\stepsTo \code{pow (n, 0)}      & \text{(clause 1 of \code{pow})}
\end{align*}

Exactly one step is invalid. \textbf{Identify which step is wrong, and
explain in one sentence why.}

\answerspace{4cm}

\begin{solution}
The \textbf{second step} is invalid. It claims that \code{1} \emph{reduces
to} \code{pow (n, 0)}, but \code{1} is already a value --- it doesn't step
to anything. While \code{1} $\eeq$ \code{pow (n, 0)} (since both evaluate
to \code{1}), equivalence does not imply reduction. Both steps should use
$\eeq$ throughout, since the theorem itself is an equivalence claim, not
a reduction claim.
\end{solution}

\end{document}
